\documentclass[]{article}

\usepackage{amsmath}
\usepackage{multirow}
\usepackage{natbib}
\usepackage{graphicx} 


\begin{document}

Our methodology is designed to identify weak lensing maps originating from anomalous simulation physics by modeling the probability distribution of their structural features, conditioned on known physical parameters. The framework consists of three main stages: (1) feature extraction using the Wavelet Scattering Transform (WST), (2) estimation of physical parameters from these features, and (3) calculation of an anomaly score using a Conditional Normalizing Flow.

\subsection{Dataset and preprocessing}
The primary dataset consists of 25,856 simulated weak lensing convergence maps for training and validation, and a separate test set of 10,000 maps. Each map is generated based on a unique combination of five physical parameters: the total matter density ($\Omega_m$), the amplitude of matter fluctuations ($S_8$), and three nuisance parameters describing baryonic feedback and observational effects ($T_{\rm AGN}$, $f_0$, $\Delta z$). To ensure our model is robust to observational noise, we added Gaussian noise with a standard deviation of $\sigma \approx 0.02582$ to each training map, matching the noise level of the test set.

\subsection{Feature extraction with the wavelet scattering transform}
The core of our feature extraction process is the Wavelet Scattering Transform (WST), chosen for its ability to capture the multi-scale, non-Gaussian morphological information that encodes the signatures of baryonic physics, while remaining stable to small deformations and robust to noise. We process each 2D convergence map using the \texttt{kymatio} library. The WST is configured with a maximum scale of $J=3$ and $L=8$ angular orientations for the wavelets at each scale. This cascade of wavelet convolutions and non-linear modulus operators produces a set of scattering coefficient maps.

To obtain a compact, fixed-size representation for each map, we apply a global average pooling operation across the spatial dimensions of the resulting coefficient maps. This procedure compresses each weak lensing map into a 217-dimensional feature vector, $\mathbf{x}_{\text{WST}}$. Finally, the entire set of training features is standardized by applying a Z-score normalization, and the mean and standard deviation derived from the training set are stored to normalize the test set features during inference.

\subsection{Conditional density modeling and anomaly scoring}
Our anomaly detection framework relies on accurately modeling the conditional probability density $p(\mathbf{x}_{\text{WST}} | \theta)$, where $\theta = (\Omega_m, S_8, T_{\rm AGN}, f_0, \Delta z)$ is the vector of five physical parameters. A high anomaly score is assigned to a map whose features $\mathbf{x}_{\text{WST}}$ are improbable given its estimated parameters $\hat{\theta}$. The inference process involves two neural network models.

\subsubsection{Parameter regression}
To obtain the conditioning vector $\hat{\theta}$ for an unlabeled map, we first train a Multi-Layer Perceptron (MLP) to regress the five physical parameters directly from the 217-dimensional normalized WST features. The network architecture consists of an input layer, two hidden layers with 512 and 256 neurons respectively, and a 5-neuron output layer. The MLP is trained to minimize the Mean Squared Error (MSE) between its predictions and the ground-truth parameters of the training set. This network provides a point estimate, $\hat{\theta}$, of the physical state of a given map.

\subsubsection{Conditional normalizing flow}
The core density estimation is performed by a Conditional Normalizing Flow (CNF). We employ a Masked Autoregressive Flow (MAF) composed of 5 autoregressive transforms, each with two hidden layers of 256 features. The CNF is trained to learn the complex, non-linear mapping from a simple base distribution (a standard Gaussian) to the distribution of WST features, conditioned on the ground-truth parameter vector $\theta$. The model is optimized by maximizing the conditional log-likelihood of the training data.

During inference, the final Out-of-Distribution (OoD) score for a given map is calculated as its negative log-likelihood (NLL) under the trained CNF, conditioned on the parameters $\hat{\theta}$ predicted by the MLP:
\begin{equation}
    \text{OoD Score} = - \log p(\mathbf{x}_{\text{WST}} | \hat{\theta})
\end{equation}
A high score indicates that the map's structural features are highly unlikely to have been generated by the fiducial simulation model, even after accounting for variations in the five known physical parameters.

\subsection{Evaluation metrics}
The performance of our OoD detection framework is evaluated using the partial Area Under the Curve (AUC) of the Receiver Operating Characteristic (ROC) curve. Specifically, we measure the mean True Positive Rate (TPR) in the highly restrictive False Positive Rate (FPR) regime of $[0.001, 0.05]$. This metric heavily penalizes methods that generate false alarms and rewards high-confidence detections. To estimate this metric during development, we constructed a synthetic validation set. The in-distribution (negative) class consisted of maps from 20 held-out cosmological models. The out-of-distribution (positive) class was created by multiplying the noiseless signal in these maps by a factor of 1.3 before adding noise, simulating a significant deviation in baryonic effects.

\end{document}
                